\scanpage{79}
que $H^1(\mathbb{F}_p,A^{\text{fixe}})=0$ (Lang) et % 1
$H^1(\mathbb{F}_p,M)=\text{fini}$ [trivial ; Attention qu'il s'agit % 2
de la coh. \uline{étale}, i.e la coh. \uline{galoisienne} de % 3
$\widehat{\mathbb{Z}}$ !], donc la dernière flèche verticale % 4
du diagramme $(1)$, est aussi bijective : % 5
\[
H^1(S,A)(\deleted{\unsure{p}}\ell)\xrightarrow{\sim}H^1(\bar S,\bar A)^\pi(\ell)
 \qquad\text{\uline{surj à noyau fini} ---} % 6
\]
Donc la \uline{conjecture} précédente signifie aussi % 7
\[
\boxed{H^1(\bar S,\bar A)(\ell)^\pi=\text{groupe fini}} % 8
\]

\deleted{Il y a une façon de \unsure{le|la} prouver,} % 9

C'est \inserted{aussi} $\pm$ équiv. (via des jacobiennes) à % 10
ceci : $X\longrightarrow S$ morphisme propre [lisse % 11
de dim relative $1$], $S$ régulier de % 12
type fini sur $\mathbb{Z}$. Prouver que % 13
$H^2(S,G_m)\longrightarrow H^{\unsure{2|1}}(X,G_m)$ est surjectif mod % 14
groupes finis ?? % 15
\margin{$\nwarrow$ \tikz[baseline=-.3ex]{\draw (0,0) circle (.035) (0,.03)--(.30,.52);}} % slanted

\margin{donc les seuls \unsure{inv.|invar.} qui restent sont $H^0$ et $H^2$ pour $A$, et \deleted{$\unsure{H^0}$}, $H^1$, $H^2$ pour $T_\ell(A)$} % slanted
\uline{NB} On devrait avoir $H^{\deleted{\unsure{2}}i}(S,A)(\ell)$ fini aussi % 16
pour \uuline{$i\ge3$}, $S$ de type fini sur $\mathbb{Z}$, $\ell$ premier aux % 17
car. résiduelles \unsure{de $S$|à ----|-----} (\inserted{NB} si $i\ge2$, $T_\ell(H^i(S,A)(\ell))\simeq H^i(S,T_\ell(A))$) % 18
% archive sheet number: 78
